\documentclass[oneside,openany,headings=optiontotoc,11pt,numbers=noenddot]{scrreprt}

\usepackage[a4paper]{geometry}
\usepackage[utf8]{inputenc}
\usepackage[T1]{fontenc}
\usepackage{lmodern}
\usepackage[ngerman]{babel}
\usepackage{ngerman}

\usepackage[onehalfspacing]{setspace}

\usepackage{fancyhdr}
\usepackage{fancybox}

\usepackage{rotating}
\usepackage{varwidth}

%Struktogramme
\usepackage[german,curves]{struktex}

\usepackage{pdflscape}
\usepackage{changepage}
\usepackage{graphicx}
\usepackage[bottom]{footmisc}
\usepackage{transparent}
\usepackage{graphbox}
\graphicspath{
	{Pics/PDFs/}
	{Pics/JPGs/}
	{Pics/PNGs/}
}
\usepackage{caption}
\usepackage{wrapfig}
\usepackage{marginnote}
\usepackage{tabularx}
\usepackage{dashrule}
\usepackage{soulutf8}
\usepackage{hhline}
%arydshln suppresses vertical lines in table
%\usepackage{arydshln}
\usepackage{multirow}
\usepackage{enumerate}
\usepackage[hidelinks]{hyperref}
\usepackage{listings}

\usepackage[table]{xcolor}
\usepackage{array}
\usepackage{enumitem,amssymb,amsmath}
\usepackage{interval}
\usepackage{cancel}
\usepackage{stmaryrd}
\usepackage{wasysym}
\usepackage{polynom}
\usepackage{diagbox}
\usepackage{dashrule}
\usepackage{framed}
\usepackage{mdframed}
\usepackage{karnaugh-map}
\usepackage{pdfpages}

\usepackage{blindtext}

\usepackage{eso-pic}

\usepackage{amssymb}
\usepackage{eurosym}

\usepackage[pages=some]{background}
\pagestyle{headings}
\renewcommand{\headrulewidth}{0.2pt}
\renewcommand{\footrulewidth}{0.2pt}
\newcommand*{\underdownarrow}[2]{\ensuremath{\underset{\overset{\Big\downarrow}{#2}}{#1}}}
\setlength{\fboxsep}{5pt}
\newcommand{\explainBelow}[3]{\underbrace{#1}_{\parbox{\widthof{#3}}{\footnotesize\raggedright #2}}}
\newcommand{\explainAbove}[3]{\overbrace{#1}^{\parbox{\widthof{#3}}{\footnotesize\raggedright #2}}}
\newcommand\footnoteref[1]{\protected@xdef\@thefnmark{\ref{#1}}\@footnotemark}


% Codestyle defined
\definecolor{codegreen}{rgb}{0,0.6,0}
\definecolor{codegray}{rgb}{0.5,0.5,0.5}
\definecolor{codepurple}{rgb}{0.58,0,0.82}
\definecolor{backcolour}{rgb}{0.95,0.95,0.92}
\definecolor{deepgreen}{rgb}{0,0.5,0}
\definecolor{darkblue}{rgb}{0,0,0.65}
\definecolor{mauve}{rgb}{0.40, 0.19,0.28}
\colorlet{exceptioncolour}{yellow!50!red}
\colorlet{commandcolour}{blue!60!black}
\colorlet{numpycolour}{blue!60!green}
\colorlet{specmethodcolour}{violet}

%Neue Spaltendefinition
\newcolumntype{L}[1]{>{\raggedright\let\newline\\\arraybackslash\hspace{0pt}}m{#1}}
\newcolumntype{M}{>{\centering\arraybackslash}X}
\newcommand{\cmnt}[1]{\ignorespaces}
%Textausrichtung ändern
\newcommand\tabrotate[1]{\rotatebox{90}{\raggedright#1\hspace{\tabcolsep}}}

%Intervall-Konfig
\intervalconfig {
	soft open fences
}

%Bash
\lstdefinestyle{BashInputStyle}{
	language=bash,
	basicstyle=\small\sffamily,
	backgroundcolor=\color{backcolour},
	columns=fullflexible,
	backgroundcolor=\color{backcolour},
	breaklines=true,
}
%Java
\lstdefinestyle{JavaInputStyle}{
	language=Java,
	backgroundcolor=\color{backcolour},
	aboveskip=1mm,
	belowskip=1mm,
	showstringspaces=false,
	columns=flexible,
	basicstyle={\footnotesize\ttfamily},
	numberstyle={\tiny},
	numbers=none,
	keywordstyle=\color{purple},,
	commentstyle=\color{deepgreen},
	stringstyle=\color{blue},
	emph={out},
	emphstyle=\color{darkblue},
	emph={[2]rand},
	emphstyle=[2]\color{specmethodcolour},
	breaklines=true,
	breakatwhitespace=true,
	tabsize=2,
}
%Python
\lstdefinestyle{PythonInputStyle}{
	language=Python,
	alsoletter={1234567890},
	aboveskip=1ex,
	basicstyle=\footnotesize,
	breaklines=true,
	breakatwhitespace= true,
	backgroundcolor=\color{backcolour},
	commentstyle=\color{red},
	otherkeywords={\ , \}, \{, \&,\|},
	emph={and,break,class,continue,def,yield,del,elif,else,%
		except,exec,finally,for,from,global,if,import,in,%
		lambda,not,or,pass,print,raise,return,try,while,assert},
	emphstyle=\color{exceptioncolour},
	emph={[2]True,False,None,min},
	emphstyle=[2]\color{specmethodcolour},
	emph={[3]object,type,isinstance,copy,deepcopy,zip,enumerate,reversed,list,len,dict,tuple,xrange,append,execfile,real,imag,reduce,str,repr},
	emphstyle=[3]\color{commandcolour},
	emph={[4]ode, fsolve, sqrt, exp, sin, cos, arccos, pi,  array, norm, solve, dot, arange, , isscalar, max, sum, flatten, shape, reshape, find, any, all, abs, plot, linspace, legend, quad, polyval,polyfit, hstack, concatenate,vstack,column_stack,empty,zeros,ones,rand,vander,grid,pcolor,eig,eigs,eigvals,svd,qr,tan,det,logspace,roll,mean,cumsum,cumprod,diff,vectorize,lstsq,cla,eye,xlabel,ylabel,squeeze},
	emphstyle=[4]\color{numpycolour},
	emph={[5]__init__,__add__,__mul__,__div__,__sub__,__call__,__getitem__,__setitem__,__eq__,__ne__,__nonzero__,__rmul__,__radd__,__repr__,__str__,__get__,__truediv__,__pow__,__name__,__future__,__all__},
	emphstyle=[5]\color{specmethodcolour},
	emph={[6]assert,range,yield},
	emphstyle=[6]\color{specmethodcolour}\bfseries,
	emph={[7]Exception,NameError,IndexError,SyntaxError,TypeError,ValueError,OverflowError,ZeroDivisionError,KeyboardInterrupt},
	emphstyle=[7]\color{specmethodcolour}\bfseries,
	emph={[8]taster,send,sendMail,capture,check,noMsg,go,move,switch,humTem,ventilate,buzz},
	emphstyle=[8]\color{blue},
	keywordstyle=\color{blue}\bfseries,
	rulecolor=\color{black!40},
	showstringspaces=false,
	stringstyle=\color{deepgreen}
}

\lstset{literate=%
	{Ö}{{\"O}}1
	{Ä}{{\"A}}1
	{Ü}{{\"U}}1
	{ß}{{\ss}}1
	{ü}{{\"u}}1
	{ä}{{\"a}}1
	{ö}{{\"o}}1
}

% Neue Klassenarbeits-Umgebung
\newenvironment{worksheet}[3]
% Begin-Bereich
{
	\newpage
	\sffamily
	\setcounter{page}{1}
	\ClearShipoutPicture
	\AddToShipoutPicture{
		\put(55,761){{
				\mbox{\parbox{385\unitlength}{\tiny \color{codegray}BBS I Mainz, #1 \newline #2
						\newline #3
					}
				}
			}
		}
		\put(455,761){{
				\mbox{\hspace{0.3cm}\includegraphics[width=0.2\textwidth]{../../logo.pdf}}
			}
		}
	}
}
% End-Bereich
{
	\clearpage
	\ClearShipoutPicture
}

\geometry{left=2.50cm,right=2.50cm,top=3.00cm,bottom=1.00cm,includeheadfoot}

\begin{document}
	\begin{worksheet}{BGY 16}{Klassenstufe 13}{Lernabschnitt 1: Exponentialfunktion - Übungen}
				
		\noindent
		\sffamily
		\begin{framed}
			\noindent
			\section*{Bevölkerungswachstum}
			Im Jahr 1996 lebten \(5,8 Mrd.\) Menschen auf der Erde. Seit dieser Zeit hat sich die Weltbevölkerung um \(2\%\) pro Jahr erhöht.
			\begin{itemize}
				\item Wie viele Menschen leben demnach heute auf der Erde?
				\item Wie viele könnten es im Jahr 2050 sein?
			\end{itemize}
		\end{framed}
		\begin{framed}
			\section*{Wir falten...}
			Ein paar Höhen, Längen und Entfernungen.\\
			\par\noindent
			\begin{tabularx}{\textwidth}{lX}
				Mainzer Dom (östliche Seitentürme) & \(55 m\)\\
				Länge der Titanic & \(269 m\)\\
				Frankfurter Fernsehturm & \(337 m\)\\
				Empire State Building & \(381 m\)\\
				Feldberg & \(1.277 m\)\\
				Mount Everest & \(8.848 m\)\\
				Entfernung zum Mond & \(399.043 km\)
			\end{tabularx}\\
			\par\noindent
			Ein handelsübliches Blatt Papier hat eine Dicke von knapp \(0,128 mm\).\\
			Angenommen Sie könnten dieses Blatt Papier ohne Einschränkung so oft halbieren, wie sie wünschen.
			\begin{itemize}
				\item Überlegen Sie, wie dick das gefaltete Papier nach einer unterschiedlichen Anzahl von Faltungen ist.\\
				\textit{Halten Sie diese Überlegungen in einer Tabelle fest.}
				\item Stellen Sie Vermutungen auf, wie häufig Sie das Papier falten müssten, und die verschiedenen Höhen, Längen oder Entfernungen zu erreichen.
				\item Welche Funktion würde die Dicke des Papiers in Abhängigkeit der Anzahl der Faltungen widerspiegeln?
			\end{itemize}
		\end{framed}
		\newpage
		\begin{framed}
			\section*{Geldanlage}
			\begin{framed}
				\noindent
				Die 12\%-Bank lockt mit folgendem Angebot:\\
				\begin{center}
					\colorbox{green!15}{\Large{12\% Zinsen nach einem Jahr!}}\normalsize
				\end{center}
				\par\bigskip\noindent
				Klingt das nicht nach einem interessanten Angebot?\\
			\end{framed}
			\begin{framed}
				\noindent
				Die benachbarte 6\%-Bank lockt hingegen mit dem folgendem Angebot:\\
				\begin{center}
					\colorbox{green!15}{\Large{6\% Zinsen nach einem halben Jahr!}}\normalsize
				\end{center}
				\par\bigskip\noindent
				Klingt doch nach einem noch interessanteren Angebot?\\
			\end{framed}
			\begin{framed}
				\noindent
				Die 3\%-Bank möchte ihre Konkurrenten ausstechen und bietet folgendes Angebot:\\
				\begin{center}
					\colorbox{green!15}{\Large{3\% Zinsen nach drei Monaten!}}\normalsize
				\end{center}
				\par\bigskip\noindent
				Gutes Angebot denken Sie bestimmt!\\
			\end{framed}
			\begin{framed}
				\noindent
				Die 1\%-Bank hat von den Angeboten der Konkurrenz Wind bekommen und muss nun nachziehen. Sie veröffentlichen das folgende Angebot:\\
				\begin{center}
					\colorbox{green!15}{\Large{1\% Zinsen nach einem Monat!}}\normalsize
				\end{center}
				\par\bigskip\noindent
				Klingt ja noch besser!\\
			\end{framed}
			Auf welches der Angebote würden Sie im ersten Moment eingehen?\\
			\par\bigskip\noindent
			Vermutung: \rule{0.8\textwidth}{0.1pt}\\
			\par\bigskip\noindent
			Überprüfen Sie ihre Vermutung und beantworten dabei die Frage: Wie viel Geld stünde ihnen nach \textbf{5 Jahren} zur Verfügung, wenn Sie zu Beginn 100\euro{} anlegen?\\
			\par
			\textit{Legen Sie eine Tabelle mit den Spalten: \underline{Jahr, Zinsen und Kapital} an. Befüllen Sie diese!}
		\end{framed}
		\newpage
		\begin{framed}
			\noindent
			\section*{Schilddrüsenuntersuchung}
			Zur Untersuchung von Schilddrüsenerkrankung wird häufig die so genannte Szintigraphie benutzt. Dabei wird radioaktives Jod-123 gespritzt, das sich in der Schilddrüse anreichert.\\
			Durch die radioaktive Strahlung kann ein Bild von der Schilddrüse aufgenommen werden.\\
			\par\bigskip\noindent
			Das radioaktive Jod-123 hat eine Halbwertszeit von ca. 13 Stunden, das heißt, das sich die Menge des Jods-123 im Laufe einer Halbwertszeit auf die Hälfte der Anfangsmenge reduziert.\\
			\begin{itemize}
				\item Wieviel Jod-123 ist nach einer Woche, verglichen mit dem Anfangswert, noch im Körper?
			\end{itemize}
		\end{framed}
	\end{worksheet}
\end{document}